\section{列表互异性导致的障碍}\label{sec:obstruction}

固定一个有限域 \(k\)。考虑这样一个检测：它接受一个列表当且仅当其中的记录两两不同。排序提供了一种实施该检测的方法。我们的下界适用于计算该检测的每个张量网络，包括用一个网络同时检测许多列表的情形。矛盾来自对每个输入比特分摊成本的估计。如果矩阵乘法指数低于 \(9/4\)，那么经过固定次数的压缩，这一成本就会降至下界以下。

\subsection{比特网络的两个估计}

一个元数为 \(u\) 的\emph{比特张量}是一个函数 \(\{0,1\}^u\to k\)。在张量网络中，每个顶点带有一张这样的表。一条连线将两个参数位置配对，每个未配对的位置是一个带名称的外部比特。给定外部比特的赋值，将各顶点相应的表项相乘，再对内部连线的所有赋值求和。这样就得到该网络所表示的张量。若每个顶点的元数至多为三，就称网络是\emph{三价的}。网络的\emph{规模}是顶点数与外部比特数之和。网络的\emph{体积}是顶点数与全部连线数之和，其中包括外部连线。由三价性，这两个数至多相差一个绝对常数因子。对于元数不受限制的网络，我们计算顶点数与所有顶点的元数之和。

第一个估计从任意一个三价网络出发，其结论为一个张量幂给出了一个网络，每份副本都有各自的外部比特。在这些比特上，新网络必须给出各份原答案的乘积。

\begin{theorem}[张量幂的压缩]\label{thm:shaving}
假设 \(\omega_k<9/4\)。则对每个 \(\eta>0\)，存在 \(\gamma>0\)，使得下述结论对于所有充分大的 \(n\) 成立。令 \(h=\lceil\log_2 n\rceil\)。若张量 \(V\) 由一个规模为 \(m\le n^{1.05}\)、具有 \(e\ge n\) 个外部比特的三价 \(k\)-网络表示，那么存在整数 \(d\ge1\)，使得张量幂 \(V^{\otimes d}\) 可用一个具有 \(de\) 个不同外部比特的三价 \(k\)-网络表示，且这一网络的规模至多为
\begin{equation}\label{eq:shave}
                   d\big(mh^{-\gamma}+eh^\eta\big).
\end{equation}
这里，固定 \(k,\eta\) 后，\(d\le n^{o(1)}\) 对输入网络一致成立。
\end{theorem}

第~\ref{sec:shaving} 节证明定理~\ref{thm:shaving}。对于每份副本，第一项使原规模除以 \(h\) 的一个正幂，第二项对每个外部比特计入 \(h^\eta\) 的成本。适当选择 \(\eta\) 便可将这一成本与列表互异性检测的成本相比较。

对 \(w\ge1\)，令 \(D_w\) 为对应于一个列表的零一张量，其中列表含有 \(2^w\) 个字，每个字长 \(2w\) 比特。张量取值一当且仅当这些字两两不同。它有 \(n_w=2w\,2^w\) 个外部比特。张量 \(D_w^{\otimes t}\) 分别检测 \(t\) 个独立给定的列表中各字的互异性。网络允许跨列表连接顶点。用 \(\operatorname{vol}_k\) 表示在 \(k\) 上实现张量所需的最小三价网络体积。

\begin{theorem}[列表互异性检测的体积]\label{thm:many}
存在仅依赖于 \(k\) 的常数 \(\sigma>0\)，使得对所有充分大的 \(w\) 以及每个整数 \(1\le t\le 2^w\)，都有
\begin{equation}\label{eq:many}
                  \operatorname{vol}_k(D_w^{\otimes t})
                          >t n_w w^\sigma .
\end{equation}
\end{theorem}

第~\ref{sec:list} 节证明定理~\ref{thm:many}。由于定理~\ref{thm:shaving} 中的复制因子由压缩过程选取，这里关于 \(t\) 的一致性至关重要。

\subsection{在固定输入长度下作比较}

\begin{proposition}\label{prop:finitefield}
在任意有限域 \(k\) 上，都有 \(\omega_k\ge 9/4\)。
\end{proposition}
\begin{proof}
对 \(2^w\) 个记录排序的一个排序网络使用 \(O(2^w w^2)\) 次比较~\citep{batcher}。比较两个 \(2w\) 比特的字并按条件交换，需要 \(O(w)\) 个布尔门。排序后，对所有相邻的字互不相等这一组条件取逻辑合取，即可算出 \(D_w\)。为用张量网络实现此电路，对每个门使用其零一表，在扇出处用相等关系的表复制比特，并将输出固定为一。对每个被接受的输入，内部连线恰有一种相容的赋值。由此得到一个定义在 \(k\) 上、规模 \(m_0=O(n_w w^2)\) 的三价网络。

假设 \(\omega_k<9/4\)。选取 \(0<\eta<\sigma\)，从定理~\ref{thm:shaving} 取出相应的 \(\gamma\)，并固定一个整数 \(J>2/\gamma\)。置 \(n=n_w\)，\(h=\lceil\log_2 n\rceil=w+O(\log w)\)。在定理的全部 \(J\) 次应用中，这些参数保持不变。在第 \(j\) 步，记列表数为 \(t_j\)，网络规模为 \(m_j\)，每个外部比特分摊的规模为 \(x_j=m_j/(t_jn)\)。初始时 \(t_0=1\)，\(x_0=O(w^2)\)。令复制因子为 \(d_j\)，由式~\eqref{eq:shave} 得
\[
         t_{j+1}=d_jt_j,\qquad
         x_{j+1}\le h^{-\gamma}x_j+h^\eta .
\]
归纳可知，在前 \(J\) 步的每一步，网络都至少有 \(n\) 个外部比特，并且 \(t_j=n^{o(1)}\)，因为 \(d_j\) 的界一致成立且 \(J\) 固定。由递推式得 \(x_j=O(w^2+h^\eta)\)，故当 \(w\) 充分大时，\(m_j=t_j n x_j\le n^{1.05}\)。这就验证了每一步均满足规模假设。对递推式求和得
\[
       x_J \le h^{-J\gamma}x_0+
                     h^\eta\sum_{r=0}^{J-1}h^{-r\gamma}
               = O(w^\eta).
\]
由于 \(\log_2 t_J=o(w)\)，定理~\ref{thm:many} 可用于这 \(t_J\) 个列表。由三价性，其网络体积为 \(O(t_J n w^\eta)\)，当 \(w\) 充分大时与式~\eqref{eq:many} 矛盾。
\end{proof}

\subsection{从有限域推广到任意域}

实现特化的关键步骤，是抽取一个在秩上取得严格改进的双线性恒等式。其系数生成一个有限生成环，其方程模掉适当理想后在有限域中仍成立。

\begin{lemma}[形式电路的双线性化]\label{lem:bilinearize}
设 \(F\) 上一个计算两个 \(s\times s\) 矩阵之积的形式电路使用了 \(q\) 个非标量乘法门或除法门，则
\begin{equation}\label{eq:gate-rank}
                     \operatorname{rank}_{\overline F}(T_s)\le 2q .
\end{equation}
这里，一个乘法门计入 \(q\) 当且仅当其两个输入均非常数，一个除法门计入 \(q\) 当且仅当其除数非常数。
\end{lemma}
\begin{proof}
在无限域 \(\overline F\) 上，选取矩阵 \(A,B\)，使得所有中间有理函数在此处均正则，且每个除数均非零。在 \(u^2=v^2=0\) 的条件下，将 \(A+uX,B+vY\) 输入电路。每个门的值有唯一的展开式
\[
       f=f_{00}+u f_{10}+v f_{01}+uv f_{11}.
\]
其中 \(f_{00}\) 为标量，\(f_{10}\) 关于 \(X\) 线性，\(f_{01}\) 关于 \(Y\) 线性。对于乘法门 \(f=ab\)，混合项的系数满足
\[
        f_{11}=a_{00}b_{11}+a_{11}b_{00}
                         +a_{10}b_{01}+a_{01}b_{10}.
\]
只有最后两项需要新的线性形式乘积。对于除法门 \(f=a/b\)，利用 \(bf=a\) 与 \(b_{00}\ne0\) 可得
\[
        b_{00} f_{11}=a_{11}-b_{11}f_{00}-b_{10}f_{01}-b_{01}f_{10}.
\]
这也至多需要两个新的乘积。线性运算，包括乘以或除以域中的标量，均不产生新的乘积。各输出的混合项系数正是 \(XY\) 的各元素。这些系数均位于至多 \(2q\) 个乘积张成的空间中，每个乘积由一个关于 \(X\) 的线性形式与一个关于 \(Y\) 的线性形式相乘而得，从而得到式~\eqref{eq:gate-rank}。
\end{proof}

\begin{lemma}[特化到有限域]\label{lem:specialize}
若在某个域 \(F\) 上有 \(\omega_F<9/4\)，则在某个有限域 \(k\) 上也有 \(\omega_k<9/4\)。
\end{lemma}
\begin{proof}
对某个固定的 \(\epsilon>0\)，电路规模的界为 \(O(s^{9/4-\epsilon})\)。由引理~\ref{lem:bilinearize}，在充分大的阶数 \(s\ge2\) 处，可取 \(T_s\) 在 \(\overline F\) 上的一个秩分解，其项数满足 \(r<s^{9/4}\)。令 \(R\subseteq\overline F\) 为在 \(\mathbb Z\) 上由该分解的所有系数生成的含幺环。分解恒等式是 \(R\) 上的多项式方程，目标张量的系数取零或一。

选取 \(R\) 的一个极大理想。其剩余域 \(k\) 作为 \(\mathbb Z\)-代数是有限生成的。在正特征下，由 Zariski 引理，\(k\) 是其素域的有限扩张。在零特征下，对 \(\mathbb Q\) 应用同一引理，则 \(k\) 是数域。用同一个正整数 \(M\) 乘以其各个代数生成元，就能将它们变成代数整数。因此，原生成元所生成的环将包含在 \(k\) 的整数环对 \(M\) 作可逆化所得的环中。对于每个素数 \(\ell\nmid M\)，该环都不包含 \(1/\ell\)，而域 \(k\) 却包含整个 \(\mathbb Q\)。因此 \(k\) 必为有限域。将上述恒等式模此极大理想约化，得 \(\operatorname{rank}_k(T_s)\le r\)。

最后，对这个固定的双线性算法取 \(j\) 次张量幂。三个线性映射中的每一个都可逐个张量坐标施行。由展平秩可知 \(s^2\le r\)，所以每个中间数组至多含有 \(r^j\) 个元素。在阶数为 \(s^j\) 时，线性运算与 \(r^j\) 次标量乘法的总成本为 \(O_{r,s}(j r^j)\) 次运算。将一般阶数填充至下一个 \(s\) 的幂，便得到 \(\omega_k\le\log_s r<9/4\)。
\end{proof}

\begin{proof}[定理~\ref{thm:square} 的证明]
命题~\ref{prop:finitefield} 和引理~\ref{lem:specialize} 给出了任意域上的下界。任意域上的上界由 Amanov、van den Berg 和 Zuiddam 证明~\citep{abvz}。
\end{proof}
